\chapter{偶数余量为什么需要多照顾一个素数}\label{app:lower}
\section{先看清楚奇偶分支的来源}
写 $r=2t,m=p+t$，把两块对称地写成
\[A=tu+v,qquad B=tu-v,qquad C=-pu.
\]
两个大块的差是 $2v$，第一块与载体的差是 $mu+v$。所以本原配置的 $G$ 是
\[G=\gcd(2v,mu+v).
\]
若 $t$ 奇，则 $m$ 偶，$\gcd(v,m)=1$ 已保证 $v$ 奇，也保证所有危险差异。若 $t$ 偶，则 $m$ 奇，除了 $\gcd(v,m)=1$，还要 $u+v$ 奇。额外的模二条件就这样出现了。

因此这里的分类不是因为喜欢分奇偶，而是同一个差分公式在素数二处真的有不同表现。后面“二进深度”只是在进一步问：偶的那个坐标到底被多少个二整除？
\InsertProof{prop:even}

\section{用一个严格下降的小步骤摆脱高二进深度}
在 $(a,z)$ 坐标里，合法性使两者奇偶相反。若偶数坐标被 $2^\nu$ 整除、$\nu\ge2$，希望一次完整返回后，两个输出共同含一个二，但不共同含四。约去这一个二，本原参数的深度就下降一。

这是很熟悉的思路：如果第一输出为“被四整除的一项”加“恰被二整除的一项”，那么它恰被二整除。我们只需要选择返回矩阵的某个系数也恰含一个二。下面的配比就是为这个小目标设计的。
\InsertProof{lem:binary-depth}

\section{两个相邻反射先造出一个足够深的群}
标准反射由两个不交平均组组成，再留下 $r$ 份旧值作载体。参数 $s$ 表示第一组取几份旧载体；不同的可行 $s$ 给不同反射。

取 $s=r-1$ 和 $s=r-3$，得到的平方标量相差 $n$，而两者互素。比较它们，就取得一个三角倍率；再与换块比较，取得剪切。这和奇数余量的证明做的是同一件事，只是系数不同。

开始得到的根步长含一个额外的 $r$ 因子。设 $t=2^a t_0$，其中 $t_0$ 奇，接下来需要把 $t_0$ 变成环的单位。难点在于它可能还整除当前控制模数，不能直接用“行列式安全”引理。证明因此设计一个夹乘 $C_cJC_c/c$，使它与 $J$ 行列式相同；除以 $J$ 后就是行列式一矩阵，可以取得逆。

\section{读完整证明时，分清三条路}
当 $t$ 奇时，局部化掉奇部以后已经足以接共同终端，路线很短。当 $t$ 偶且有三个标准反射时，要先调整二进方向，再调整奇素数方向，最后匹配共同单位。两个边界只有两个标准反射，则加一个迹为 $\pm2$ 的返回，直接把根步长改善到 $2n$，又能接共同终端。

第二条路最长。读它时可以把“方向”理解为比值 $z/a$，把“尺度”理解为同一方向的哪个本原列。先把比值调对，再把共同倍数调对，两步缺一不可。
\InsertProof{lem:even-core}

\section{乘三、变号，为什么能处理所有二进单位}
证明里有 $\kappa\equiv3\pmod8$。模八的奇数只有 $1,3,5,7$，乘三和变号就能从一到达全部四类。更高位怎么办？

最简单地看 $3^2=9=1+8$。平方一次，$9^2-1=80$ 恰被 $16$ 整除；再平方，赋值继续增加一。一般地，若 $u=1+2^bc$、$b\ge2,c$ 奇，则
\[u^2-1=2^{b+1}c+2^{2b}c^2
\]
恰含 $b+1$ 个二。这样逐位选择幂，就能修正每个更高二进位。所谓“生成主单位核”，在这里就是这件具体的逐位调整。

\section{r接近p时，重新选择实际分组}
当 $r=p-1$，要保留的载体太大，刚才某些标准分组不再有足够位置。于是先平均一组制造新值，再把其中 $p-1$ 份留作载体，之后用第二组及补集恢复两个 $p$ 块。

它仍是相邻反射、倍率、剪切的同一机制，只是必须重新核对取用数。最小的 $p=5,n=14$ 再单列三个具体矩阵；这些矩阵的行列式因子和夹乘等式都列在下面。
\InsertProof{lem:upper-boundary}
这一讲可以看出，一般代数工具已经很统一，长度主要花在“现有几份数够不够实现那个工具”。位置容量不能由群论自动补出来。
